\clearpage

\section{SYSTEM DESIGN}

System design translates the analytical requirements and clinical preprocessing pipelines into a robust, deployable machine learning architecture. This section presents the formal design, modular implementation, algorithmic mechanics, and data-partitioning protocols designed to classify unstructured medical transcription narratives into clinical specialties.

\subsection{MODEL(S) BUILDING}

The classification framework employs a multi-model paradigm comprising four distinct supervised learning families: linear margin optimization (Support Vector Machine), probabilistic link modeling (Logistic Regression), non-parametric bagging (Random Forest), and Bayesian conditional independence (Multinomial Naive Bayes). Furthermore, an Ensemble Voting meta-classifier is synthesized to consolidate individual algorithmic strengths into an optimal consensus predictor.

\subsubsection{Implementation Code}

The model training and inference pipeline is implemented in Python using \texttt{scikit-learn} and \texttt{imbalanced-learn}. The system executes sublinear TF-IDF vectorization, cost-sensitive loss calibration, probability calibration, and ensemble voting aggregation. Figure \ref{lst:model_implementation} provides the core implementation of the model building pipeline.

\begin{figure}[H]
\centering
\begin{minipage}{0.96\linewidth}
\begin{small}
\begin{verbatim}
import joblib
import pandas as pd
import numpy as np
from sklearn.model_selection import train_test_split
from sklearn.feature_extraction.text import (
    TfidfVectorizer
)
from sklearn.linear_model import LogisticRegression
from sklearn.svm import LinearSVC
from sklearn.ensemble import (
    RandomForestClassifier,
    VotingClassifier
)
from sklearn.naive_bayes import MultinomialNB
from sklearn.calibration import CalibratedClassifierCV
from sklearn.metrics import (
    accuracy_score,
    classification_report,
    confusion_matrix
)

# 1. Load Filtered 8-Specialty Dataset
df = pd.read_csv('filtered_mtsamples_8classes.csv')
X = df['clinical_text'].astype(str)
y = df['medical_specialty']

# 2. Stratified 80/20 Train-Test Split
X_train, X_test, y_train, y_test = train_test_split(
    X, y,
    test_size=0.20,
    random_state=42,
    stratify=y
)

# 3. Controlled TF-IDF Feature Extraction
tfidf = TfidfVectorizer(
    max_features=8000,
    ngram_range=(1, 2),
    sublinear_tf=True,
    min_df=2,
    max_df=0.7,
    norm='l2'
)
X_train_tfidf = tfidf.fit_transform(X_train)
X_test_tfidf = tfidf.transform(X_test)

# 4. Base Classifiers with Class Balancing
lr = LogisticRegression(
    class_weight='balanced',
    max_iter=1000,
    random_state=42
)
svm = LinearSVC(
    class_weight='balanced',
    max_iter=3000,
    random_state=42
)
cal_svm = CalibratedClassifierCV(
    LinearSVC(class_weight='balanced', max_iter=3000)
)
rf = RandomForestClassifier(
    n_estimators=100,
    max_depth=25,
    class_weight='balanced',
    random_state=42
)
mnb = MultinomialNB(alpha=0.5)

# 5. Soft Voting Ensemble (Weighted Consensus)
voting_ensemble = VotingClassifier(
    estimators=[
        ('lr', lr),
        ('svm', cal_svm),
        ('rf', rf),
        ('mnb', mnb)
    ],
    voting='soft',
    weights=[2, 2, 1, 1],
    n_jobs=1
)

# 6. Model Training & Evaluation
voting_ensemble.fit(X_train_tfidf, y_train)
y_pred = voting_ensemble.predict(X_test_tfidf)

print("Accuracy:", accuracy_score(y_test, y_pred))
print(classification_report(y_test, y_pred))
\end{verbatim}
\end{small}
\end{minipage}
\caption{Core Implementation Code for 4-Model Ensemble Pipeline}
\label{lst:model_implementation}
\end{figure}

\subsubsection{Model Planning (Training/Testing Split of the Dataset)}

Model planning governs the experimental protocol to guarantee unbiased evaluation and eliminate data leakage. The dataset consists of 1,663 curated clinical transcriptions spanning 8 mutually exclusive clinical departments.

To preserve the underlying specialty distribution while evaluating generalizability on unseen patient records, an \textbf{80:20 Stratified Partitioning Protocol} was established:
\begin{itemize}
    \item \textbf{Training Partition (80\%)}: Comprises 1,330 clinical records used exclusively for vocabulary extraction, inverse document frequency parameter estimation, and classifier weight optimization.
    \item \textbf{Testing Partition (20\%)}: Comprises 333 clinical records completely held out during feature extraction and model fitting, reserved strictly for final inference validation.
    \item \textbf{Stratification Protocol}: The split is governed by the class distribution vector $y$, ensuring that minority specialties (such as Ophthalmology and ENT) maintain proportional representation across both partitions.
    \item \textbf{Data Leakage Prevention}: The TF-IDF vectorizer parameters ($\text{IDF}(t)$, vocabulary dictionary, document frequencies) are fit \textit{strictly} on $X_{\text{train}}$. The test partition $X_{\text{test}}$ is transformed using the fixed vocabulary parameters learned from training.
\end{itemize}

Table \ref{tab:train_test_split_breakdown} details the exact sample distribution across the training and testing sets for each of the 8 clinical specialties.

\begin{table}[H]
\centering
\small
\resizebox{\linewidth}{!}{
\setlength{\tabcolsep}{6pt}
\begin{tabular}{|l|c|c|c|c|}
\hline
\textbf{Clinical Specialty} & \textbf{Total Records} & \textbf{Train Samples (80\%)} & \textbf{Test Samples (20\%)} & \textbf{Class Proportion} \\
\hline
Cardiovascular / Pulmonary & 371 & 297 & 74 & 22.31\% \\
Orthopedic & 355 & 284 & 71 & 21.35\% \\
Gastroenterology & 224 & 179 & 45 & 13.47\% \\
Neurology & 223 & 178 & 45 & 13.41\% \\
Urology & 156 & 125 & 31 & 9.38\% \\
Obstetrics / Gynecology & 155 & 124 & 31 & 9.32\% \\
ENT - Otolaryngology & 96 & 77 & 19 & 5.77\% \\
Ophthalmology & 83 & 66 & 17 & 4.99\% \\
\hline
\textbf{Total Corpus} & \textbf{1,663} & \textbf{1,330} & \textbf{333} & \textbf{100.00\%} \\
\hline
\end{tabular}
}
\caption{Stratified Training and Testing Split Distribution across 8 Clinical Specialties}
\label{tab:train_test_split_breakdown}
\end{table}

\subsubsection{Training}

\paragraph{Build the Model Using the Train Dataset}
During the training phase, the four candidate classifiers and the meta-voting ensemble are parameterized using the 1,330 training instances vectorized across the 8,000-dimensional TF-IDF space.

\begin{enumerate}
    \item \textbf{Logistic Regression (LR) Optimization}:
    Trained as a multinomial log-linear classifier minimizing the penalized cross-entropy objective:
    \begin{equation}
        \mathcal{L}_{\text{LR}}(\mathbf{W}) = -\sum_{i=1}^{N} \sum_{k=1}^{K} w_k \cdot y_{i,k} \log P(y_i = k \mid \mathbf{x}_i; \mathbf{W}) + \frac{1}{2C} \|\mathbf{W}\|_F^2
    \end{equation}
    where $w_k = \frac{N}{K \cdot N_k}$ represents the class-balancing weight inversely proportional to class frequency $N_k$, penalizing classification errors on under-represented specialties such as Ophthalmology and ENT. Optimization is resolved via the L-BFGS quasi-Newton solver over 1,000 iterations.

    \item \textbf{Support Vector Machine (Linear SVM) Optimization}:
    Trained as an unconstrained maximum-margin formulation using one-vs-rest (OvR) decomposition. The primal optimization objective is formulated as:
    \begin{equation}
    \begin{aligned}
        \min_{\mathbf{w}_k, b_k, \boldsymbol{\xi}} \quad & \frac{1}{2} \|\mathbf{w}_k\|_2^2 + C \sum_{i=1}^{N} w_{y_i} \xi_{i,k} \\
        \text{subject to} \quad & y_{i,k}(\mathbf{w}_k^T \mathbf{x}_i + b_k) \ge 1 - \xi_{i,k}, \quad \xi_{i,k} \ge 0
    \end{aligned}
    \end{equation}
    The linear kernel is computationally optimal for sparse, high-dimensional textual matrices ($D = 8,000$), resisting empirical overfitting. To allow probability integration within the Soft Voting Ensemble, a duplicate SVM model is calibrated via Platt scaling using 5-fold cross-validation (\texttt{CalibratedClassifierCV}).

    \item \textbf{Random Forest (RF) Optimization}:
    Constructed as an ensemble of $B = 100$ bootstrap decision trees constrained to maximum depth $d_{\text{max}} = 25$. At each split node, a random subset of $m = \sqrt{8,000} \approx 89$ features is evaluated against the Gini impurity criterion:
    \begin{equation}
        I_G(p) = 1 - \sum_{k=1}^{K} p_k^2
    \end{equation}
    Balanced bootstrap subsampling (\texttt{class\_weight='balanced'}) dynamically recalculates sample weights for each bootstrap tree, preventing dominant surgical and cardiovascular notes from monopolizing the split nodes.

    \item \textbf{Multinomial Naive Bayes (MNB) Optimization}:
    Computes prior probabilities $P(C_k)$ and feature likelihood parameters $\theta_{k,j} = P(w_j \mid C_k)$ from training count matrices with Laplace smoothing parameter $\alpha = 0.5$:
    \begin{equation}
        \theta_{k,j} = \frac{N_{k,j} + \alpha}{N_k + \alpha \cdot |V|}
    \end{equation}
    where $N_{k,j}$ is the aggregate TF-IDF weight of word $j$ within specialty $k$, and $|V| = 8,000$ represents the vocabulary size.

    \item \textbf{Voting Ensemble Consensus Synthesis}:
    The predictions of the four trained estimators are pooled via two ensemble voting schemes:
    \begin{itemize}
        \item \textbf{Hard Voting}: Assigns the predicted specialty by majority vote:
        \begin{equation}
            \hat{y}_{\text{hard}} = \operatorname{argmax}_k \sum_{m=1}^{M} w_m \cdot \mathbb{I}(\hat{y}_m = k)
        \end{equation}
        \item \textbf{Soft Voting}: Averages the posterior probability distributions across models weighted by empirical confidence:
        \begin{equation}
            \hat{y}_{\text{soft}} = \operatorname{argmax}_k \sum_{m=1}^{M} w_m \cdot P_m(y = k \mid \mathbf{x})
        \end{equation}
        where weights are assigned as $w_{\text{LR}} = 2$, $w_{\text{SVM}} = 2$, $w_{\text{RF}} = 1$, and $w_{\text{MNB}} = 1$.
    \end{itemize}
\end{enumerate}

\subsubsection{Testing}

\paragraph{Evaluate the Performance of the Model Using the Test Dataset}
Testing executes a strictly separated inference pipeline on the 333 unobserved clinical narratives in the testing partition $X_{\text{test}}$:

\begin{enumerate}
    \item \textbf{Inference Pipeline Flow}:
    Each test record undergoes identical preprocessing transformations: lowercasing, punctuation stripping, stop-word elimination, and projection into the 8,000-dimensional vector space using the pre-fitted vectorizer:
    \begin{equation}
        \mathbf{x}_{\text{test}} = \operatorname{TFIDF}(t_{\text{test}}; \mathbf{V}_{\text{train}})
    \end{equation}
    \item \textbf{Independent Prediction Generation}:
    The vector is fed into each standalone model and the voting ensemble pipelines:
    \begin{equation}
        \hat{\mathbf{y}}_{\text{test}} = \mathcal{M}(\mathbf{X}_{\text{test\_tfidf}})
    \end{equation}
    \item \textbf{Metric Evaluation Protocol}:
    Predicted labels $\hat{\mathbf{y}}_{\text{test}}$ are benchmarked against true ground-truth labels $\mathbf{y}_{\text{test}}$. Performance is quantified across global exact accuracy, macro-averaged precision, macro-averaged recall, class-specific F1-scores, and weighted F1-metrics to rigorously verify robustness against remaining class skewness.
\end{enumerate}
